\section{The Lean Substrate}
\label{sec:substrate}

The grammar developed in Part~\ref{sec:grammar} is constructible.
A companion technical paper \cite{ref_chio_parent_paper} develops the substrate in detail.
This Part sketches the substrate's relevant features without re-deriving them.
The objective is to give the legal-academy reader enough familiarity with the formal machinery to assess
whether the grammar is in fact buildable in the sense the Article claims,
without reproducing the technical content that a separately published paper has already established.

The substrate is a system in which a polity is represented as a triple
$(T, C, K)$,
where $T$ is a receipt-admission scope, $C$ is a Merkle-rooted citizenship roster of the actors authorized to seek
admission of receipts into the polity's history, and $K$ is a finite list of predicates that the polity has
adopted as its constitution.
A receipt is admissible into the polity's history only if the predicates in $K$ accept it,
and the polity's history is the totality of admitted receipts.

The substrate provides four operations on this triple that are relevant to the grammar:
admission, capability attenuation, treaty intersection, and constitutional amendment.
Each operation is implemented in Rust and accompanied by Lean~4 proofs of its key properties.
The properties most relevant to this Article's grammar are the following.

\subsection{Admission as decidable predicate evaluation}

A receipt $r$ is admissible into a polity $(T, C, K)$ if and only if $r \in T$ and $K(r) = \text{accept}$.
The predicate evaluation $K(r)$ is decidable in the sense that the constitution $K$ is a finite list of
predicates and each predicate is evaluable as a Boolean function over the receipt's canonical representation.
The substrate's runtime kernel implements this evaluation as a fail-closed admission hook:
a receipt for which any predicate returns reject is denied,
and a receipt for which all predicates accept is admitted and signed.

The grammar's typed-rollback witness corresponds, on the substrate, to a predicate in $K$ that requires the receipt
to carry a rollback obligation.
A receipt for a bounded executive action is constructed with fields for $\text{act}$, $\TTL$, $w$, and $q$,
and the predicate in $K$ that gates admission of such receipts requires all four fields to be present and to type-check.
A receipt that lacks a witness $w$, or that carries a witness that does not type-check against the substrate's
witness specification,
fails the predicate evaluation and is denied.

\subsection{Capability attenuation and the executive authority}

The substrate represents authority through capability tokens that are signed objects carrying rights.
A capability for a bounded executive action carries the action's specification and is signed by the authority
that granted it.
The substrate's capability-attenuation operation permits the holder of a capability to derive a more restricted
capability from the original,
and a Lean theorem \cite{ref_chio_parent_paper} establishes that the derived capability cannot grant rights
that the original did not.

The relevance to delegated emergency authority is that the structural relationship between a legislature granting
an emergency authority to an executive officer is, on the substrate's framing, an instance of capability delegation
with attenuation.
The legislature is the original authority;
the executive officer holds an attenuated capability whose scope is bounded by the legislature's grant.
The attenuation theorem ensures that the executive officer cannot, by exercise of the delegated authority,
acquire rights the legislature did not grant.
The grammar's typed-rollback witness is an attenuation constraint at the capability level:
the executive officer's capability is bounded by the requirement that any action taken under it
carry a witness for the path back.

\subsection{Constitutional amendment and backward refinement}

The substrate's amendment operation transforms a polity from $(T, C, K)$ to $(T, C, K')$ where $K'$ is a
candidate replacement constitution.
The amendment is admissible only if $K'$ is a backward refinement of $K$:
every receipt admitted under $K$ remains admissible under $K'$.
A Lean theorem on the substrate establishes that the backward-refinement check is decidable and that the
substrate's runtime kernel refuses amendment requests that do not exhibit a refinement witness
\cite{ref_chio_parent_paper}.

The relevance to delegated emergency authority is that an emergency-amendment proposal --
a proposal to modify the polity's constitution in response to an exceptional circumstance --
is admissible only if it preserves backward refinement.
The amendment cannot, by its own operation, retroactively un-admit a receipt that the polity has previously admitted.
The grammar's typed-rollback witness operates at the action level rather than the amendment level,
but the structural intuition is the same:
admission of a receipt or amendment requires construction of a witness for compatibility with the polity's
prior history.

\subsection{The substrate's relationship to the legal claim}

The substrate is not the Article's contribution.
The substrate is cited as evidence that the grammar developed in Part~\ref{sec:grammar} is buildable.
A legal-academy reader does not need to verify the Lean proofs to assess the Article's structural claim about
delegated emergency authority.
The reader needs only to accept that the grammar is, in principle, implementable on a substrate of the kind
the companion paper describes.
The substrate's existence answers the objection that the grammar is a theoretical construct without operational meaning.
It does not answer the objection that the grammar may be inappropriate for the legal domain to which the Article applies it.
That objection is the subject of Part~\ref{sec:limits}.
